\documentclass[12pt]{article}
% Préambule commun à tous les documents extraits de « La Revue CAPES »
\usepackage[T1]{fontenc}
\usepackage[utf8]{inputenc}
\usepackage{lmodern}
\usepackage{amsmath,amssymb,amsthm,amsfonts,mathrsfs}
\usepackage{setspace}
\usepackage{listings}
\usepackage{pgfplots}
\pgfplotsset{compat=1.18}
\usepackage{longtable}
\usepackage{adjustbox}
\usepackage{lettrine}
\usepackage{tkz-tab}
\usepackage{changepage}
\usepackage{pdflscape}
\usepackage{graphicx}
\usepackage{tikz}
\usepackage{xcolor}
\usepackage[a4paper,top=2cm,bottom=2.5cm,left=2.5cm,right=2cm]{geometry}
\usepackage{fancyhdr}
\usepackage{draftwatermark}
\SetWatermarkText{La RevueCapes}
\SetWatermarkLightness{0.9}
\SetWatermarkScale{2}
\usepackage{hyperref}
\hypersetup{colorlinks=true,linkcolor=blue!50!black,urlcolor=blue!50!black}

\definecolor{revue}{RGB}{20,60,120}

\pagestyle{fancy}
\fancyhf{}
\renewcommand{\headrulewidth}{0.4pt}
\setlength{\headheight}{15pt}

% \entete{Matière}{Titre}{Type de document}
\newcommand{\entete}[3]{%
  \fancyhead[L]{\small\textsc{La Revue CAPES} -- #1}%
  \fancyhead[R]{\small #2}%
  \fancyfoot[C]{\thepage}%
  \thispagestyle{plain}%
  \begin{center}
    {\color{revue}\rule{\textwidth}{1.2pt}}\\[4pt]
    {\small\textsc{Concours directs CAP-CEG / CAPES -- Burkina Faso}}\\[6pt]
    {\LARGE\bfseries\color{revue} #2}\\[6pt]
    {\large\itshape #1 -- #3}\\[2pt]
    {\color{revue}\rule{\textwidth}{1.2pt}}
  \end{center}
  \vspace{0.5cm}%
}

% Encadré pour les remarques ajoutées dans les corrigés
\newcommand{\remarque}[1]{\par\medskip\noindent\fcolorbox{revue}{revue!6}{\parbox{\dimexpr\linewidth-2\fboxsep-2\fboxrule}{\textbf{\color{revue}Remarque.} #1}}\par\medskip}


\begin{document}
\entete{Mathématiques}{Sujet type n°3}{Énoncé}
\begin{spacing}{1.5}

 
 \textbf{Exercice 1 }
 
  Soit $f :\mathbb{R}^2\longrightarrow \mathbb{R}$ la fonction ainsi définie par:
 $f(x,y)=2x^2+2y^2+2xy-x-y$
 
 1. Calculer les dérivées partielles premières de $f$.
 
 2. Déterminer l'unique point critique $A$ de $f$.
 
 3. Calculer les dérivées partielles secondes de $f$.
 
 4.montrer que $f$ présente en $A$ un extremum en précisant sa nature.
 
 \textbf{Exercice 2}
 
  On considère les déterminants suivants:
 $\Delta_1= \begin{vmatrix}
 1 & 1 & 1\\
 b & c & d\\
 b^2 & c^2 & d^2
 \end{vmatrix}$
 et $\Delta_2=\begin{vmatrix}
 1 & 1 & 1 & 1 \\
 a & b & c & d \\
 a^2 & b^2 & c^2 & d^2 \\
 a^3 & b^3 & c^3 & d^3
\end{vmatrix} $
 
 
 1. Calculer $\Delta_1$ en donnant le résultat sous forme factorisée.
 
 2. Montrer que $\Delta_2 =(b-a)(c-a)(d-a)\Delta_1$.
 
 3. En déduire $\Delta_2$ sous forme factorisée.
 
 4. En déduire le rang de $\Delta_2$ suivant les valeurs de $a$ et $b$.
 
 \medskip\textbf{Exercice 3 }
 
 On considère la fonction $f$ définie sur $\mathbb{R}$ par $f(x) = 2e^{-2x}sin(x)$.
 
 1. Déterminer le développement limité à l'ordre 3 au voisinage de 0 de $f$.
 
 2. En déduire
 \[
 \lim_{x\longrightarrow 0}\frac{2e^{-2x}sin(x)}{x}
 \]

 \textbf{Exercice 4 }
 
 On lance deux dés rouges et verts et on note $A$ l'évènement "La somme des numéros fait 6" et
 $B$ l'évènement "Sur le dé rouge, on obtient un nombre pair".
 
 1. Définir ce que sont deux évènements incompatibles.

 2. Les évènements $A$ et $B$ sont-ils incompatibles ?
 
 3. Quand dit-on que deux évènements sont indépendants?
 
 4. Calculer $P(A)$ et $P(B)$.
 
 5. Calculer $P(A\cap B)$ et $P(A\bigtriangleup B)$.
 
 6. Que peut-on déduire sur les deux évènements $A$ et $B$?
 
 
 \quad
 \quad
 

\end{spacing}
\end{document}
